\chapter*{Appendix 14: Scenario Xi (Architecture)}

\section*{1. The Legacy Dependency (Technical \\\& Cultural)}
Physical architecture is heavily constrained by state power: municipal zoning laws, building codes, legacy real estate debt (mortgages), and fiat property taxes. The Sovereign Stack requires physical space for nodes, communal living, and production, which currently exists on territory claimed by the state.

\section*{2. The Parasitic Bridge (Technical \\\& Legal)}
Layer 7 employs aggressive "Legal Parasitism." Sovereign physical spaces are purchased or leased using cooperative LLCs or religious/educational trusts, exploiting legacy tax loopholes and zoning exceptions. The architecture itself is modular and non-permanent (e.g., mobile nodes, tiny homes), technically classifying as "temporary structures" to evade the most restrictive legacy building codes while tapping into the physical grid where necessary.

\section*{3. The Decoupling Trigger}
The legal decoupling is triggered when the sovereign network reaches a physical and economic density that makes state enforcement of zoning laws practically impossible or economically ruinous for the municipality. At this point, the physical spaces transition from "legally compliant under loopholes" to "openly sovereign autonomous zones."

\section*{4. Orchestration \\\& Ethical Mechanics}
The orchestration requires Layer 5 spatial mapping to optimize the placement of sovereign infrastructure within hostile legacy territory. Ethically, this process must avoid gentrifying vulnerable neighborhoods; instead, it should focus on reclaiming abandoned commercial real estate or utilizing the bridge to provide superior infrastructural services (mesh internet, solar power) to the surrounding legacy community, gaining local cultural immunity.
